% Generated by build_paper_tex.py from stabcert.md. Do not edit by hand.
% Never typeset by its generator -- no LaTeX toolchain was available.
\documentclass[11pt]{article}
\usepackage[utf8]{inputenc}
\usepackage[T1]{fontenc}
\usepackage{amsmath,amssymb,amsthm}
\usepackage{booktabs}
\usepackage{hyperref}
\usepackage[margin=1in]{geometry}

\newtheorem{theorem}{Theorem}
\newtheorem{lemma}{Lemma}
\newtheorem{proposition}{Proposition}
\theoremstyle{definition}
\newtheorem{definition}{Definition}
\newtheorem{hypothesis}{Hypothesis}

\title{StabCert: translation validation for stabilizer channels}
\author{Frédéric Kramarczyk\\Independent researcher\\\texttt{fkra62@gmail.com}}
\date{}

\begin{document}
\maketitle




\begin{abstract}

Quantum routing compilers rewrite circuits to satisfy hardware connectivity
constraints, but the resulting circuits are seldom verified against the channel
they are meant to implement. We present a decision procedure for the equality
of stabilizer channels --- Clifford unitaries with stabilizer ancilla preparation
and partial trace --- restricted to a specified input subspace.

Our procedure compares the canonical signed signatures of code-Choi states,
relative to a support determined by the specification, and is both sound and
complete: two such channels agree on the code subspace if and only if their
signatures coincide. It runs in polynomial time, avoids dense matrix
representations and basis state enumeration, and its verdict is invariant under
any change of Stinespring dilation.

We implement this in StabCert, a tool that treats the compiler as an untrusted
black box and certifies routed circuits produced by Qiskit SABRE and pytket
against targets reconstructed from a specification rather than from a reference
circuit. On the families studied, our measurements identify generator density ---
not circuit width alone --- as a key factor in verification cost.


\end{abstract}
\section{Introduction}

Quantum programs written for hardware must be routed: a compiler inserts SWAP
gates and relabels qubits so that every two-qubit operation acts on an edge of
the device coupling map. This transformation is substantial, and it is
performed by heuristic search procedures that are not themselves verified.
Qiskit's SABRE and pytket's routing passes are widely deployed and largely
trusted on the strength of testing.

Existing verification approaches address adjacent problems.

Compiler certification --- proving the transformation itself, once and for all ---
applies only where the compiler's source is available and under the verifier's
control. SABRE and pytket are neither.

Equivalence checking, as implemented for example in QCEC, compares a compiled
circuit against a reference circuit. Like our approach it relies on canonical
representations; the distinction is not canonicity but the object represented
and the input required. Section 8 states it precisely, together with the
capabilities and limits of the engines involved.

Deductive verification proves properties of programs written in a dedicated
language, reaching algorithms as complex as Shor's order finding, but takes its
input as a program rather than as an artifact produced elsewhere.

None of these answers the question a routing pass raises in practice: given a
circuit produced by an untrusted compiler, does it implement the channel the
specification calls for? The distinction matters because the reference circuit
is itself an artifact. Comparing against it certifies faithfulness to a
possibly wrong starting point.

\textbf{Contribution.} We restrict attention to stabilizer channels, as defined in
Section 2, and show that within this class the question is decidable in
polynomial time. Specifically:

\begin{itemize}
\item We define the code-Choi state, a mixed stabilizer state --- a state stabilized
by a signed subgroup --- encoding a channel's action on a specified input
subspace, and prove that equality of its canonical signed signature is
necessary and sufficient for equality of the channels on that subspace
(Section 3).
\item We show the verdict is invariant under any change of Stinespring dilation, so
that circuits differing only in ancilla convention are correctly accepted
(Section 3.4).
\item We implement the procedure in StabCert and certify routed circuits produced
by SABRE and pytket, with adversarial campaigns reporting no false accepts
and no false rejects (Sections 5--6).
\item We measure verification cost and separate the contribution of circuit width
from that of generator density (Section 7).
\end{itemize}

The property is proved of the decision procedure; the adversarial campaigns of
Section 6 test the implementation that realises it, and no finite corpus could
establish the former.

\textbf{Scope.} The class is narrow by design. Gottesman--Knill makes stabilizer
circuits classically simulable, which is what permits an exact decision
procedure; it also means our results say nothing about non-Clifford gates,
noise, or timing. Section 4 states the boundary precisely.


\section{Preliminaries}

We assume familiarity with the stabilizer formalism and fix notation.

An $n$-qubit \textbf{stabilizer group} is an abelian subgroup of the Pauli group not
containing $-I$. It is represented by a \textbf{tableau}: for each generator, a
binary vector $(x|z) \in F_2^{2n}$ together with a sign. Group operations
correspond to $F_2$ row operations on the binary parts, with signs carried
along. The word \emph{tableau} is reserved for this representation throughout.

A \textbf{Clifford unitary} maps Pauli operators to Pauli operators under
conjugation, and is fully specified by the images of $2n$ generators. A
\textbf{stabilizer channel} is a Clifford unitary composed with stabilizer ancilla
preparation and partial trace over a discarded register.

By the \textbf{Gottesman--Knill theorem}, stabilizer circuits admit efficient
classical simulation: a tableau of $2n$ signed generators determines the state,
and each gate updates it in polynomial time. This is what makes the decision
procedure of Section 3 possible, and it also bounds it --- everything below is
confined to the stabilizer fragment.

\subsection{Registers and symbols}

\begin{center}
\begin{tabular}{llllll}
\toprule
symbol & object \\
\midrule
$R$ & purifying reference for the message, ` & M & ` wires \\
$M$ & message, channel input \\
$X$ & accessible partition, channel output \\
$C$ & inaccessible complement, traced \\
$O \subseteq X$ & requested output, with ` & O & = & M & ` \\
$\Lambda$, $\Lambda'$ & stabilizer channels under comparison \\
$N$, $N^c$ & the channel from $M$ to $X$, and its complement to $C$ \\
$S$ & stabilizer group of the source Choi state \\
$S_X$ & subgroup of $S$ surviving the partial trace onto $X$ \\
$\tau_X$ & marginal $N(I/d)$ on $X$, with `d = 2^{ & M & }` \\
$\Pi$ & projector onto $supp(\tau_X)$ \\
$k$ & logical qubits, `k = & X & - & S_X & ` \\
$E$ & support encoder (Definition~\ref{def:1}) \\
$Ref(X)$ & reference register of the code subspace, ` & X & ` wires \\
$sig(\cdot)$ & canonical signed signature (Definition~\ref{def:3}) \\
\bottomrule
\end{tabular}
\end{center}

$R$ and $Ref(X)$ are distinct registers of different widths and are never
interchanged.


\section{Deciding equality of stabilizer channels}

\subsection{Setting}

Fix a channel specification $P$, consisting of an input register $M$, a set of
ancilla qubits with initial stabilizers, a source Clifford $U$, and an
accessible partition $X$.

Let $S$ be the stabilizer group generated by the specification's input
generators conjugated by $U$, and let

\begin{verbatim}
S_X = { S|_X : S ∈ S, supp(S) ⊆ X }
\end{verbatim}

be the subgroup surviving the partial trace onto $X$. Write
$\Pi = \prod_i (I + S_i)/2$ for the projector onto the joint $+1$ eigenspace of $S_X$,
and $k = \lvert X \rvert - \lvert S_X \rvert$ for the number of logical qubits. The target subspace is
$supp(\tau_X)$ with $\tau_X = \Pi / 2^k$.

\begin{hypothesis}[specification-determined support]\label{hyp:1}
$\Pi$ is a function of $P$
alone. No part of the candidate artifact contributes to its construction.
\end{hypothesis}

This is what the verdict below is relative to, and it is the hypothesis on
which the theorem's usefulness rests: a candidate may claim a support, but that
claim is compared against $\Pi$, never used to build it. $\Pi$ is never
materialised; it is represented throughout by its signed generators.

\subsection{The code-Choi state}

\begin{definition}[support encoder]\label{def:1}
Let $\bar{X}_i, \bar{Z}_i$ ($i < k$) be canonical
logical representatives of the code defined by $S_X$, and $D_j, S_j$ its
destabilizer--stabilizer pairs. The support encoder $E$ is the Clifford on
$\lvert X \rvert$ qubits determined by
\begin{equation*}\begin{split}
E X_i E^\dagger = \bar{X}_i ,      E Z_i E^\dagger = \bar{Z}_i        (i < k) \\
E X_{k+j} E^\dagger = D_j ,  E Z_{k+j} E^\dagger = S_j    (j \ge 0) \\
\end{split}\end{equation*}
\end{definition}

Both families of labels are obtained by canonical $F_2$ elimination, so $E$ is
determined by $P$ and involves no gauge choice. $E$ maps
$\lvert \psi\rangle_k \otimes  \rvert0\rangle^{\lvert X \rvert-k}$ into the code space, and carries the signs on the
stabilizer tail.

The problem format requires $O \subseteq X$ with $\lvert O \rvert = \lvert M \rvert$, enforced at problem
construction; this is the condition under which the two sides of the comparison
below are of the same type.

\begin{definition}[code-Choi state]\label{def:2}
Let $Ref(X)$ be a reference register of
$\lvert X \rvert$ qubits. Prepare $k$ Bell pairs between $Ref(X)$ and the first $k$ wires
of $X$, with the remaining $\lvert X \rvert - k$ wires of each half in $|0\rangle$. Apply the
\textbf{complex conjugate} $Ē$ to $Ref(X)$ and the support encoder $E$ to $X$. Pass
$X$ through the candidate channel $\Lambda$. The code-Choi state $J_\Pi(\Lambda)$ is the
result reduced onto $O \cup Ref(X)$.
\end{definition}

The conjugation on the reference is not a convention of convenience. It is the
same fact as the transpose appearing in the inversion of Theorem~\ref{thm:1}, seen from
the other side: $Ē$ on $Ref(X)$ is what makes $Tr_Ref[(\sigma^T \otimes I) J_\Pi(\Lambda)]$
return $\Lambdã(\sigma)$ rather than $\Lambdã(\sigma^T)$. Stating either without the other leaves
the definition and the proof inconsistent with each other.

Three remarks, each of which a reader will otherwise have to reconstruct.

\emph{The reduction registers differ on the two sides.} The target is reduced onto
$R \cup X$; the candidate onto $O \cup Ref(X)$. They are comparable because
$\lvert O \rvert = \lvert M \rvert$, which Definition~\ref{def:1} requires.

\emph{Three distinct partial traces occur in the chain, and they must not be
conflated.} $R$ and $C$ are traced to form $\tau_X$; $C$ is traced to form the
reduced target; $X \setminus O$ is traced to form the reduced candidate. Only the last
belongs to Definition~\ref{def:2}.

\emph{The state is mixed, on $\lvert O \rvert + \lvert X \rvert$ wires.} It is not a purification and not a
pure stabilizer state: tracing $X \setminus O$ leaves a state stabilized by a proper
subgroup. What the construction avoids is a product of projectors --- it returns
an independent generating set directly, so no normalisation scalar propagates
through the canonical form.

The full reference register is kept, not truncated to $k$ wires, because the
encoder mixes the $\lvert X \rvert - k$ padding wires into the code. After encoding no
factor is separable, and no subset can be discarded.

\subsection{Canonical form and the decision procedure}

\begin{definition}[canonical signed signature]\label{def:3}
For a set of independent signed
Pauli generators, put the binary vectors $(x|z)$ in reduced row echelon form,
performing each row operation simultaneously on the binary vector and on the
corresponding Pauli string, so that phases follow the elimination rather than
being recomputed. The canonical signed signature $sig(\cdot)$ is the resulting
ordered tuple of signed generators.
\end{definition}

\begin{lemma}\label{lem:1}
\emph{Let the generators be independent and generate a group not
containing $-I$. Then the canonical signed signature is a normal form: two such
sets generate the same signed subgroup if and only if their signatures are
equal.}
\end{lemma}

\begin{proof}
RREF over $F_2$ is unique for a given row space, so the binary parts
agree exactly when the subgroups agree as unsigned groups. Because each row
operation is applied to the Pauli strings in parallel, the phase attached to
each echelon generator is the phase of the unique group element with that
binary support --- unique precisely because $-I \notin S$, without which the group
would contain both $P$ and $-P$ for the same support. Hence equality of
signatures is equivalent to equality of the signed subgroups.
\end{proof}

The hypothesis $-I \notin S$ is not incidental: the implementation rejects any
specification whose generators fail it, together with commutation and
independence, before any comparison is attempted.

\begin{theorem}[soundness and completeness]\label{thm:1}
\emph{Let $\Lambda$, $\Lambda'$ be stabilizer
channels and $\Pi$ as above. Then}
\begin{equation*}\begin{split}
sig(J_\Pi(\Lambda)) = sig(J_\Pi(\Lambda'))  \iff  \Lambda(\rho) = \Lambda'(\rho) for all \rho with supp(\rho) \subseteq supp(\tau_X) \\
\end{split}\end{equation*}
\end{theorem}

\begin{proof}
(⇐) Suppose the channels agree on the code subspace. The states $\rho$
supported in $supp(\Pi)$ span the operator space on that subspace, so by
linearity the two channels agree on all of it, hence on the joint state of
Definition~\ref{def:2}. The code-Choi states coincide, and by Lemma~\ref{lem:1} their signatures
are equal.

(⇒) Equal signatures give equal signed stabilizer groups, hence
$J_\Pi(\Lambda) = J_\Pi(\Lambda')$ as states. Three steps then recover the channels.

First, $E$ restricts to a unitary isomorphism between $C^{2^k} \otimes \lvert 0\rangle^{ \rvertX|-k}$
and the code subspace. By Definition~\ref{def:1}, $E$ maps the stabilizer group of
$\lvert 0\rangle^{ \rvertX|-k}$ on the last $\lvert X \rvert - k$ wires onto $S_X$, so it carries that
subspace onto the joint $+1$ eigenspace of $S_X$, which is $supp(\Pi)$; being
Clifford it is unitary, hence a bijection between the two. Every $\rho$ with
$supp(\rho) \subseteq supp(\tau_X)$ is therefore $E(\sigma \otimes \lvert 0\rangle\langle0 \rvert)E^\dagger$ for a unique $\sigma$ on $k$
qubits.

Second, $J_\Pi(\Lambda)$ is the ordinary Choi state of the induced logical channel
$\Lambdã : \sigma \mapsto \Lambda(E(\sigma \otimes \lvert 0\rangle\langle0 \rvert)E^\dagger)$. The prepared state of Definition~\ref{def:2} is
$(Ē \otimes E)(Φ_k \otimes \lvert 0\rangle\langle0 \rvert)(Ē \otimes E)^\dagger$, whose reduction on $X$ is $\Pi/2^k$ and whose
correlations with $Ref(X)$ are those of $Φ_k$ transported by $Ē \otimes E$. Applying
$\Lambda$ to $X$ and reducing onto $O \cup Ref(X)$ therefore yields $(id \otimes \Lambdã)(Φ_k)$ up
to the relabelling $Ē$ performs on the reference --- which is precisely the
transpose carried in the inversion below.

Third, Choi--Jamio\l{}kowski inverts explicitly:

```
Λ̃(σ) = 2^k · Tr_Ref[ (σ^T ⊗ I) · J_Π(Λ) ]
```

where the transpose on the reference is the convention the construction
applies, with $Y^T = -Y$. Equal Choi states therefore give $\Lambdã = \Lambdã'$, and by
the first step $\Lambda = \Lambda'$ on every state supported in $supp(\tau_X)$.

Equality of the induced logical channels determines $\Lambda$ only on the image of
$E$, which is $supp(\tau_X)$ --- precisely the domain the theorem asserts.
\end{proof}

The implementation compares canonical signatures. An equivalent double-inclusion
test exists in the same module but has no caller; Lemma~\ref{lem:1} makes the two
equivalent, and the signature form is retained because it is a normal form ---
one comparison rather than $2\cdot\lvert generators \rvert$ membership solves.

\subsection{Gauge invariance}

\begin{proposition}\label{prop:1}
\emph{The verdict is invariant under any transformation acting
only on the discarded environment.}
\end{proposition}

\begin{proof}
$J_\Pi$ is obtained by partial trace over the environment, and the trace
is invariant under isometries acting on the traced factor alone.
\end{proof}

Consequently two circuits differing only in ancilla convention or Stinespring
dilation receive the same verdict --- a property required for accepting output
from compilers that do not share the reference implementation's conventions.
Section 6.2 reports a family of artifacts that exercises this restriction
rather than assuming it.


\section{Scope}

The class we treat is narrow, and the boundary is a design choice rather than
an artifact of the implementation. We state it before reporting results so that
every subsequent claim is read within it.

\subsection{What the procedure covers}

Stabilizer channels, as defined in Section 2. Within this class the decision
of Section 3 is exact.

Gottesman--Knill is what makes this possible, and the same fact bounds the
result: nothing here transfers to non-Clifford gates, and the procedure says
nothing about circuits containing $T$ gates or arbitrary rotations.

\subsection{What is out of scope by construction}

Noise, calibration, real-time decoding, and scheduling. These are physical or
temporal properties; a stabilizer Choi state does not express them. Verifying
that a circuit implements the intended channel is separable from verifying that
a device executes it faithfully, and we address only the former.

\subsection{Measurement and feed-forward}

The recovery circuits verified here are coherent isometric dilations: they
contain no mid-circuit measurement, and consequently no branching on
measurement outcomes.

The class \textbf{admits} a symbolic treatment of the general case: a Pauli
correction that is an $F_2$-linear function of a syndrome is a linear map, which
a stabilizer formalism can verify once rather than per branch. The current
implementation \textbf{contains no syndrome handling at all} --- its recovery circuits
are measurement-free, so no branch arises and no such matrix is constructed. We
\textbf{claim nothing} about that path and report no results for it.

Readers familiar with the Yoshida--Kitaev decoder should note that its canonical
form is measurement-based; the circuits here are its coherent counterpart.

\subsection{What is verified and what is reported}

The verdict covers semantic equality of the reduced channel on the specified
subspace, conformance to the coupling map --- every two-qubit gate is checked
against an edge of the device graph --- and invariance under change of
Stinespring dilation.

Depth and two-qubit gate counts are recomputed by the verifier from the final
gate list and \textbf{enter the verdict}. SWAP attribution --- how many SWAPs serve
movement versus restoration --- is not reconstructible from the artifact format
and participates in no verdict; the tables of Section 6 report it as a
measurement.

\subsection{Bounds on the cost measurements}

The scaling results of Section 7 are measured on one family of instances: a
chain architecture with random stabilizer scramblers, a single seed, and two
scrambling depths. Two consequences.

First, the exponent we confirm is an exponent of that family. Because the
family's generator density drifts with $n$, it is not a degree at fixed
density; Section 7.6 separates the two.

Second, no extrapolation to surface-code circuits is warranted. The densities
measured range from \textbf{0.149 to 0.578}. For a surface code the same metric is
derived rather than measured: a code of distance $\Delta$ has $\Delta^2$ data qubits and
stabilizer generators of weight at most four, so mean weight $w̄$ gives density
$w̄ / \Delta^2$, and $\Delta = 5$ sits at $3.33 / 25 \approx 0.133$ --- at the edge of the measured
range rather than beyond it. The order-of-magnitude gap appears only at
$\Delta \ge 13$, where constant-weight generators over a growing lattice drive the
density as $1/n$. Cost projections beyond the measured range are upper bounds
on a family chosen to be unfavourable, not predictions.

\subsection{Two adversarial campaigns}

We report two campaigns with different objects, and the figures are not
interchangeable. \texttt{orelia.verifier-adversarial-campaign/v1} targets the verifier
as a whole; \texttt{orelia.channel-certified-adversarial-campaign/v1} targets the
channel-certified policy specifically. Section 6.2 states which is which; a
reader encountering either figure in a published CSV can identify it by these
format identifiers. Conflating them would suggest an inconsistency where there
is none.


\section{Implementation}

StabCert implements the procedure of Section 3. It is a Python package with two
runtime dependencies, NumPy and Stim, the latter providing the tableau engine
on which canonical comparison rests. It is available under Apache 2.0.

The package also contains a routing procedure. It exists because the
\texttt{reproducible-route} policy below requires a reference route to compare
against; it is not offered as a compiler and no claim is made about its output
relative to other routers.

\subsection{Two policies}

The tool exposes two verification policies, and the distinction is the point of
the design.

\texttt{reproducible-route} requires bit-for-bit equality with the reference route. It
answers \emph{did this run reproduce the reference implementation?} and is used for
regression testing.

\texttt{channel-certified} reconstructs the target from the channel specification and
accepts any synthesis, routing, or Stinespring dilation whose canonical signed
signature matches. It answers the question of Section 1: \emph{does this artifact,
however produced, realise the specified channel on the specified subspace?}

Circuits routed by third-party compilers fail the first policy and pass the
second. That is the intended behaviour, and Section 6 reports it as such.

\subsection{The certified closure}

The verification path is a closed import closure, computed by AST traversal
from the three entry points --- compiler, verifier, command line --- and enforced
by the test suite. The closure comprises eleven modules; the verifier alone
reaches eight. No module within it constructs a dense matrix or enumerates
basis states.

The distribution also ships exploratory and instance-construction code outside
this closure, including a dense stabilizer-group enumeration. No verification
path can reach it, and the closure test fails if that changes. We state this
because the claim of Section 3 --- polynomial time, no enumeration --- is a claim
about the closure, not about every file in the package; a reader who greps the
installed package will find the enumeration and should know why it is harmless.

The closure assertion is bidirectional: it fails both if an unreachable module
becomes reachable and if a declared module is no longer reached, so the
declaration cannot decay into a wish list.

\subsection{What the verdict aggregates}

A single invocation checks, and reports separately:

\begin{itemize}
\item Semantic equality of the reduced channel on the code subspace, by canonical
signature comparison (Section 3).
\item Coupling-map conformance: every two-qubit gate in the routed circuit is
checked against an edge of the device graph, gate by gate, on the final gate
list. No routing trace is required, since the final gate list is what the
device executes.
\item Metadata consistency: the artifact's claimed support is compared against the
reconstructed $\Pi$, never used to build it (Hypothesis~\ref{hyp:1}).
\end{itemize}

These are independent conditions, not components of a score: a failure in any
one produces a failing verdict, and no success elsewhere compensates for it.

Resource figures are recorded alongside the verdict; Section 4.4 states which
of them enter it.

\subsection{Reproducibility}

Published test counts are not written by hand. A synchronisation script replays
the suite, parses the summary, and rewrites the counted spans in the
documentation; it refuses to touch the documentation if the suite is not green,
and a \texttt{--check} mode exits non-zero on divergence for use in CI.

The script guarantees that the documentation matches the local suite. It does
not, by itself, guarantee anything about what a reader obtains from a checkout.
That is established separately: a continuous-integration job runs the suite
from a checkout on each push, under Python 3.10 and 3.12, with core
dependencies only and then with every optional backend. The suite passes: 124
tests with all backends, 104 passing and 2 modules skipped without them.


\section{Evaluation}

\subsection{Third-party routes}

We compile reference recovery circuits for three fixtures and route them with
two external compilers, Qiskit SABRE and pytket. Six routed artifacts result.
For each, the verdict is identical:

\begin{center}
\begin{tabular}{ll}
\toprule
check & result \\
\midrule
route differs from reference & yes \\
\texttt{reproducible-route} & rejected \\
\texttt{channel-certified} & accepted \\
phase mutation & rejected \\
falsified final permutation & rejected \\
deterministic replay & byte-identical artifact \\
\bottomrule
\end{tabular}
\end{center}

The first three lines are the intended behaviour of the two policies: a route
produced elsewhere is not the reference route, and is nonetheless the specified
channel. The last three establish that acceptance is not vacuous --- a single
sign flip in the tableau, or a permutation claimed but not realised, is caught.

\textbf{Reproducibility.} The routed artifacts reported here were produced with
Qiskit 2.5.1 (SABRE, $decay$ heuristic, seed 20260803, one trial) and pytket
2.18.1 (\texttt{RoutingPass} with \texttt{LexiLabellingMethod} and \texttt{LexiRouteRoutingMethod}),
on the fixtures published under \texttt{tests/fixtures/recovery\_v1/}. SABRE is a
stochastic search: a different seed or version yields a different route, which
\texttt{channel-certified} accepts and \texttt{reproducible-route} rejects --- the verdicts of
this section are stable under that variation, the resource figures are not.

\subsection{Adversarial campaigns}

Two campaigns, with distinct objects and non-interchangeable figures.

\begin{center}
\begin{tabular}{lllll}
\toprule
campaign & invalid & valid & false accepts & false rejects \\
\midrule
\texttt{orelia.verifier-adversarial-campaign/v1} & 10 000 & 1 000 & 0 & 0 \\
\texttt{orelia.channel-certified-adversarial-campaign/v1} & 1 300 & 800 & 0 & 0 \\
\bottomrule
\end{tabular}
\end{center}

One family in the second campaign deserves mention, because it is the only one
that separates the two possible specifications. Its members prefix the circuit
by an element of the stabilizer group of $\tau_X$: every state of the code
subspace is a $+1$ eigenvector, so the action there is unchanged while the
total unitary differs. These artifacts must be accepted under the
subspace-restricted specification of Section 3 and would be rejected under a
total-channel specification. \textbf{All 100 are accepted.} Without this family, no
measurement in either campaign would distinguish the two readings of the
theorem.

Mutation testing is reported by family rather than in aggregate: 10 000 sign
flips is one class tested 10 000 times, not 10 000 classes.

\subsection{Resources}

Depth and two-qubit gate counts are recomputed by the verifier from the final
gate list and enter the verdict; they are certified quantities. The \texttt{reference}
column is the route the verifier reconstructs and against which difference is
established --- it is a baseline, not a scored competitor.

\begin{center}
\begin{tabular}{llllll}
\toprule
A & arch & logical & reference & SABRE & pytket \\
\midrule
1 & chain & 12 / 14 & 46 / 62 & 49 / 68 & 34 / 41 \\
8 & chain & 82 / 111 & 391 / 687 & 584 / 939 & 401 / 694 \\
12 & grid_2d & 154 / 229 & 464 / 751 & 668 / 1177 & 454 / 808 \\
\bottomrule
\end{tabular}
\end{center}

Routing overhead, relative to the logical circuit:

\begin{center}
\begin{tabular}{lllll}
\toprule
A & arch & reference & SABRE & pytket \\
\midrule
1 & chain & 3.83 / 4.43 & 4.08 / 4.86 & 2.83 / 2.93 \\
8 & chain & 4.77 / 6.19 & 7.12 / 8.46 & 4.89 / 6.25 \\
12 & grid_2d & 3.01 / 3.28 & 4.34 / 5.14 & 2.95 / 3.53 \\
\bottomrule
\end{tabular}
\end{center}

\textbf{Both tables are to be read across, not down.} Overhead rises from $A = 1$ to
$A = 8$ and falls at $A = 12$; the fall is a topology effect, not a size effect
--- the grid offers better connectivity than the chain. Rows compare routers at a
fixed instance; they compare nothing to each other.

On $A = 12$, pytket produces a circuit shallower than the reference by 10
layers while using 57 more CNOTs. Which of the two is preferable has no answer
independent of a stated cost function --- and this is a fact about routing, not a
comparison of tools.

Between the two external routers, pytket produces both the shallower and the
cheaper circuit on all three fixtures. Three instances, two topologies, one
SABRE seed and one configuration each: this is what was observed, not a
ranking.

\subsection{Time}

\begin{center}
\begin{tabular}{ll}
\toprule
measurement & value \\
\midrule
SABRE regression, 3 fixtures, end to end & 50.06 s, peak RSS 88.95 MiB \\
pytket regression, 3 fixtures & 49.67 s, peak RSS 55.97 MiB \\
verification alone, 2 100 campaign cases & median 77.3 ms, mean 83.9 ms, max 172.1 ms \\
verification at $n = 9 / 20 / 40$ & 0.24 s / 10.43 s / 272.23 s \\
\bottomrule
\end{tabular}
\end{center}

The last row cross-checks Section 7.5 without having been fitted to it:
doubling $n$ from 20 to 40 multiplies time by \textbf{26.1}, against \textbf{27.5}
predicted by the exponent of 4.78 measured on the upper-half window over
$n = 9\dots40$.


\section{Cost}

Two different questions are answered here, and conflating them is the way this
section is most likely to be misread. \textbf{What does verification cost as
instances grow?} is a question about a family. \textbf{What is the degree at fixed
generator density?} is a question about the algorithm. The family we measure
has a density that drifts with $n$, so the two answers differ: roughly $n^5$
for the first, roughly $n^{5.7}$ for the second. Both are stated; neither is
the other.

\subsection{Counters, not seconds}

Cost is reported as GF(2) operation counts, instrumented in the elimination
routines. Three counters matter: the number of affine systems solved, the
number of row XORs, and the number of scalar bit XORs. Row XORs are the
machine-relevant measure, since NumPy vectorises a row operation; scalar bit
XORs are a machine-model-independent upper bound that overcounts vectorised
work by about $n$.

Counters are deterministic: they do not vary with machine load, scheduling, or
repetition. Section 7.5 explains why wall-clock time, which does vary, is
reported but not used as the complexity.

\subsection{An exact identity}

Over 32 instances, $n = 9$ to $40$, the number of affine systems solved is not
approximately but \textbf{exactly}

\begin{verbatim}
N_sys(n) = 28n² − 232n + 598
\end{verbatim}

with zero residual. Second differences are constant at 56; the reader can
recompute this from the published CSV. This fixes the call structure of the
nested elimination: the two nested loops compose to $n^2$ solves, not $n^3$.

\subsection{A prediction, registered before measurement}

Composing the nesting with the cost of one elimination --- $n^2$ solves, each a
Gaussian elimination costing about $n^3$ row operations, each row operation
touching about $n$ bits --- predicts degrees 2, 5 and 6 for the three counters.

Observed exponents cannot be compared to those targets directly: they are
inflated by lower-order terms. The counter whose degree is known exactly
calibrates that inflation on the same instances. The inflation is derivable in
closed form from the identity above --- the tangent exponent is $n\cdotN_sys'(n)/N_sys(n)$,
so the bias is

\begin{verbatim}
β(n) = (232n − 1196) / (28n² − 232n + 598)  →  232/28 = 8.286 / n
\end{verbatim}

matching the measured secant bias to within $1.2 \times 10^-^3$ across the range.

With $\tau = max(observed(S) - 2, 0.1)$ and $excess = observed - target$, the rule
registered before the sweep was: \textbf{confirmed} if $excess \le \tau$ and the exponent
is still falling; \textbf{rejected} if $excess > 2\tau$ or the exponent rises;
\textbf{undecided} otherwise, which alone justifies extending the range. The floor
of $0.1$ prevents the test from degenerating as the bias vanishes; it did not
bind here, the measured bias being 0.2896.

Registering this mattered. At $n \le 30$ the rule returned \emph{undecided} for
\texttt{row\_xors} by 0.007 --- an excess of 0.3914 against a tolerance of 0.3843. A
threshold chosen after seeing that number would have been chosen differently.

\subsection{Result}

Windows are the lower and upper halves of $n = 9\dots40$, 16 points each.

\begin{center}
\begin{tabular}{llllllll}
\toprule
counter & target & lower & upper & trend & excess & tolerance & verdict \\
\midrule
\texttt{affine\_systems\_solved} & 2 & 2.690 & 2.290 & -0.400 & --- & --- & calibrator \\
\texttt{row\_xors} & 5 & 6.285 & \textbf{4.907} & -1.378 & -0.093 & 0.290 & confirmed \\
\texttt{scalar\_bit\_xors} & 6 & 7.643 & \textbf{6.106} & -1.536 & +0.106 & 0.290 & confirmed \\
\texttt{verify\_seconds} & --- & 4.711 & 4.780 & \textbf{+0.068} & --- & --- & not a verdict \\
\bottomrule
\end{tabular}
\end{center}

The structural composition holds at the registered threshold. \texttt{row\_xors}
crossing below 5 is not evidence of a degree below 5: it is a data-dependent
counter that fell by 1.38 over the range, and an overshoot of 0.09 is within
that movement. Single-step exponents are not reported --- at the top of the range
they are unstable enough to be meaningless.

\subsection{Wall time is not the complexity}

Timing gives an exponent of \textbf{4.780} on the upper-half window over
$n = 9\dots40$. An unrelated instance family --- four structural preflights over
$choi_qubits = 26\dots32$ --- gives a log-log exponent of 4.55 for its
construction-and-verification time. The agreement is a consistency check across
harnesses, not a second measurement of the same quantity: the ranges, the
parameters and the timed sections differ.

Time grows \textbf{more slowly} than the operation counts (4.780 against 4.907 for
row XORs), and its exponent \emph{rises} with $n$ --- 4.711, then 4.780 --- while every
counter's falls. The cause is vectorisation: NumPy performs a row XOR in
roughly constant time, so measured seconds track row operations rather than bit
operations, and converge on that exponent from below.

The two series therefore bracket the answer from opposite sides: counters
descending from 6.285 to 4.907, time ascending from 4.711 to 4.780, meeting
within 0.13. \textbf{Two independent series converging from above and below constrain
more than a two-parameter fit to either one.}

\subsection{Density, and what the exponent is an exponent of}

The measured family's generator density is not constant: it drifts as
$n^{-0.238}$. This is mechanical rather than incidental --- at fixed scrambler
depth the number of two-qubit gates grows linearly while the number of matrix
entries grows quadratically, so normalised density must fall.

A paired design was registered to neutralise density by holding scrambler depth
fixed. It failed on its own terms: the density contrast between the two depths
itself drifts from 1.259 to 2.455, an exponent of $n^{0.582}$ against a
tolerance of 0.384, so the ratio exponent conflates a degree effect with a
widening gap and is not readable. Density does not stay neutralised; it must be
modelled.

Pooling both depths at 22 matched widths --- where the elimination structure is
identical, the same $N_sys(n)$ at every width --- gives

\begin{verbatim}
log(row_xors) = a + b·log n + c·log(density)
b = 6.313     c = 1.410
\end{verbatim}

Before reading those coefficients, two gates. The model must reproduce each
arm's observed exponent: it does, to $\pm0.077$ against a tolerance of 0.384
(dense, observed 6.054 against 5.977; sparse, observed 5.079 against 5.156).
And the same model fitted to the exactly-known counter, whose true density
elasticity is zero, must return approximately zero: it returns $c_0 = 0.006$.
The method invents essentially nothing.

$c = 1.410$ against a threshold of 0.384 and a false-positive floor of 0.006:
\textbf{density governs a real share of the cost.} $b = 6.313$ is an all-points fit
and carries the same finite-size inflation the calibrator shows on that fit
($b_0 = 2.611$ against a true 2, so $+0.611$), giving a bias-corrected degree at
fixed density of about \textbf{5.7} --- higher than the family exponent of 5, as it
must be, since density falls with $n$ and cost rises with density. The relation
that reconciles the two is $family exponent = b + c\cdot\delta$, with $\delta$ the density
drift of the arm.

One weakness to report: the maximum log residual is 0.408 for \texttt{row\_xors}
against 0.104 for the calibrator. The model satisfies the exponent gate but
fits appreciably less well pointwise; $c = 1.41$ has the right sign and order
of magnitude, and its second decimal does not.

This is why Section 4.5 forbids extrapolation to a surface code. The relevant
comparison is of drift regimes: the dense arm drifts as $n^{-0.238}$, the
sparse arm as $n^{-0.820}$, a surface code as $n^{-1}$ --- constant-weight
generators over a growing lattice. The sparse arm already approaches the
surface-code density regime without having its locality. What is captured is
density; what is not is locality.

\subsection{One method tried and rejected}

Fitting $e(n) = \gamma + K/n$ to the sequence of local exponents extrapolates the
asymptotic degree directly and would have terminated the range question
cleanly. Run first against the calibrator, whose degree is exactly 2, it
returns $\gamma = 1.941$ with a jackknife spread of $[1.939, 1.942]$ over
$n = 9\dots40$ --- excluding the truth inside an interval some thirty times too
narrow.

The $1/n$ form is only the leading correction, so its residuals are systematic,
and a jackknife measures the stability of a fit rather than its accuracy. The
method is not used. We record it because the failure is not specific to these
data: any resampling estimate of uncertainty will tighten around the wrong
value when the error is systematic, and only a quantity with a known answer
reveals it.


\section{Related work}

Three lines of work address adjacent problems. We have not executed any of the
tools discussed below; the characterisations that follow are drawn from their
published documentation and papers.

\textbf{Compiler certification.} Proving the transformation itself, once and for
all, is stronger where it applies --- and it applies only to compilers whose
source is available and under the verifier's control. SABRE and pytket are
neither. Translation validation, certifying each output rather than the
producer, is the classical answer to that situation, and it is the one we take.

\textbf{Equivalence checking.} QCEC compares a compiled circuit against a reference
circuit. Its decision-diagram engines provide exact equivalence checking by
comparing canonical representations of the implemented operators, with support
for ancilla and garbage qubits and a notion of partial equivalence defined on
measurement distributions, although resource consumption can be exponential in
the worst case; a ZX-calculus engine is often effective at establishing
equivalence but, being based on an incomplete rewriting procedure, cannot in
general establish inequivalence. The two engines are combined deliberately, and
exploiting the $G G'^\dagger$ structure is the published contribution that keeps the
diagrams small in practice.

Like our approach, QCEC relies on canonical representations. The distinction is
not canonicity but the object represented and the input required: QCEC
canonically represents operators and compares two of them, whereas we
canonically represent a stabilizer channel restricted to a code subspace and
compare one against a target reconstructed from a channel specification. This
is not a deficiency of QCEC --- in its own use case, checking that a compiler
preserved the source circuit, the reference \emph{is} the specification and the
question is well posed. The gap matters only where no reference exists yet.

\textbf{Deductive verification.} Qbricks verifies circuit-building quantum programs,
reaching parametric implementations of Shor's order finding, quantum phase
estimation and Grover's search, with high proof automation over parametrized
path sums in Why3. Its scope is far wider than ours in the class of circuits
reached; its documented workflow is to write a program in its own DSL, and we
found no documented import path for an artifact produced elsewhere. The two
tools therefore answer different questions and admit no meaningful quantitative
comparison: proof effort and verification time are not the same quantity, and
tabulating them side by side would suggest a commensurability that does not
exist.

\textbf{Summary.} The axis that separates this work from all three is the input.
Compiler certification takes a compiler; equivalence checking takes two
circuits; deductive verification takes a program. We take an artifact of
unknown provenance and a channel specification.

\bibliographystyle{plain}
\bibliography{stabcert}
\end{document}
